\section{Input for Standard Mesh} \label{sec2.2}
%

\medskip

{\bf General Remarks}

\medskip

EIRENE follows atoms, molecules or test ions in a 3-dimensional
computational box. This box is discretized by
zones (mesh cells), the boundaries of which are
defined either by regular ``standard mesh surfaces", i.e., co-ordinate
surfaces, (BREP representation of geometry, \ref{sec1.5}) which are described here, and/or by zones whose
boundaries are defined more generally
by ``additional surfaces" (see below block 3B).
Alternative to this BREP option there is also the option of unstructured meshes (Voxel representation, loc.cit.) of triangles ( in 2D)
or tetrahedra (in 3D).

In 2-dimensional applications the co-ordinate surfaces automatically collapse to
coordinate lines, and in 1-dimensional applications they collapse to coordinate (grid) points.

As long as test particles travel inside the ``standard mesh", all
cell indexing computations are done automatically. At transition into
the more general ``additional mesh" and inside these additional
cells, this indexing has
to be specified by the flags ILSWCH, ILCELL etc. by the user, in
input block 3B.

There are up to 3 sets of standard co-ordinate surfaces, each set being parameterised by the arrays
RSURF(i1), PSURF(i2) and TSURF(i3).

Here RSURF(1), RSURF(2), ...,
RSURF(NR1ST) are grid-points (parameters) in the 1st co-ordinate direction, mostly the radial or x direction
(depending on the
geometry level, internal parameter LEVGEO, see below) defined by the input data in sub-block 2A.
I.e., there are NR1STM= NR1ST-1 cells in the 1st co-ordinate direction (e.g. radial or x direction).
With regard to the terminology in \ref{sec1.5}, equation (\ref{eq1.34}), a
corresponding co-ordinate surface is labelled with a double index 1,IR
and is defined by the equation

$f_1^{1,IR} (x,y,z)$
=constant=RSURF(IR).

PSURF(1), PSURF(2), ..., PSURF(NP2ND) are grid-points in the 2nd co-ordinate direction,
typically the poloidal or y direction (sub-block 2B), and there are
NP2NDM=NP2ND-1 cells in this direction.
A corresponding
co-ordinate surface is labelled $2,IP$ and is defined by the equation

$f_1^{2,IP} (x,y,z) $=constant=PSURF(IP).

In case NP2ND = 1 this 2rd grid is absent and the calculation automatically collapses to a 1D case,
i.e. a case with only 1 set of co-ordinate points.  This does not mean that the 2rd, ignorable coordinate must be
straight cartesian (y-direction). E.g. a spherically symmetric volume, but with only radial resolution in space
and polar and azimuthal symmetry,
then is still referred to as 1-dimensional calculation.

In case of the NLPLG option (mesh in x-y plane generated by a
set of polygons)
for simplicity and by an abuse of language we sometimes refer to
the polygons themselves as radial surfaces or ``radial polygons"
RSURF and refer to the
polygons obtained by connecting all points with a particular index along
the various radial polygons as ``poloidal polygons" PSURF.
No use is made in EIRENE of the arrays RSURF and PSURF in connection
with the polygon mesh option. The NLPLG option permits to run EIRENE
on 2D computer generated meshes with grid cuts in one direction.

Furthermore, TSURF(1), TSURF(2), ...,
TSURF(NT3RD) are grid-points in the third co-ordinate direction, mostly the toroidal or z direction (sub-block 2C).
A co-ordinate surface is labeled $3,IT$ and is defined by the equation
$f_1^{3,IT} (x,y,z)$
=constant=TSURF(IT).

In case NT3RD = 1 this 3rd grid is absent and the calculation automatically collapses to a 2D case,
i.e. a case with only 2 sets of co-ordinate lines.  This does not mean that the 3rd, ignorable coordinate must be
straight cartesian (z-direction). E.g. a full torus, but with only radial and poloidal resolution and toroidal symmetry
then is referred to as 2-dimensional calculation.

In addition to these grids and, depending on the geometry options
chosen, EIRENE then computes further ``derived" grid arrays, for example
a ``(radial) surface labeling grid" RHOSRF, which may or may not be
identical to the RSURF grid,
and a ``zone centered (radial) surface labeling grid"
RHOZNE(IR) defined by:

\texttt{RHOZNE(IR) = 0.5 $\cdot$ (RHOSRF(IR)+RHOSRF(IR+1))}

These grids may be used for the input of closed form
radial- or x direction
plasma profiles or as abscissae in plots of e.g.
poloidally (or y) and toroidally
(or z) averaged tallies.

We have found it convenient to allow for several (NBMLT)
identical copies
of such ``standard meshes", separated and bounded by arbitrary
``additional surfaces".
Input flags for this ``mesh multiplication"
option are comprised in sub-block 2D.

As mentioned above,
in addition to these regular grids one can define cells of almost
arbitrary complexity by using the ``additional surfaces" defined in
sub-block 3B. Such surfaces are labeled $0,IA$  in EIRENE.
Appropriate cell number switching flags must be set
on transparent additional
surfaces separating these general ``additional cells".
Data relevant for the
additional cells (e.g. the volume)
are read in sub-block 2E. The number of additional cells
is NRADD.

Therefore, the total number of cells in an EIRENE run is NSBOX, with

NSBOX = (NR1ST $\cdot$ NP2ND $\cdot$ NT3RD) $\cdot$ NBMLT + NRADD .

Cell volumes in the standard mesh region, as well as the center of
mass vector in each of these cells are computed automatically. The
corresponding data in the additional cell region must be specified
by the user, e.g. in the input block 2E or 8, or in the user
supplied routine GEOUSR (section \ref{sec3.6}), which is called by
EIRENE in the initialization phase.

\medskip

{\bf The Input Block}

\medskip

{\sl *** 2. Data for Standard Mesh}
\begin{verbatim}
 (INDGRD(J), J=1,3)
 ** 2A. x- or radial co-ordinate surfaces
 NLRAD
 IF (NLRAD) then
    NLSLB  NLCRC  NLELL  NLTRI  NLPLG  NLFEM  NLTET NLGEN
    NR1ST  NRSEP  NRPLG  NPPLG  NRKNOT NCOOR
    IF (INDGRD(1).LE.5) THEN
       IF (NLSLB.OR.NLCRC.OR.NLELL.OR.NLTRI) THEN
          RIA  RGA  RAA  RRA
          IF (NLELL.OR.NLTRI) THEN
             EPIN  EPOT  EPCH  EXEP
             ELIN  ELOT  ELCH  EXEL
             IF (NLTRI) THEN
               TRIN  TROT  TRCH  EXTR
            ENDIF
          ENDIF
       ELSEIF (NLPLG) THEN
          XPCOR  YPCOR  ZPCOR  PLREFL
         (NPOINT(1,K) NPOINT(2,K), K=1,NPPLG)
          DO 21 I=1,NR1ST
             (XPOL(I,J) YPOL(I,J),  J=1,NRPLG)
 21       CONTINUE
       ELSEIF (NLFEM) THEN
         XPCOR  YPCOR  ZPCOR
         NRKNOT
         (XTRIAN(I),I=1,NRKNOT)
         (YTRIAN(I),I=1,NRKNOT)
         DO ITRI=1,NR1ST
           I,NVERT(1,ITRI),NVERT(2,ITRI),NVERT(3,ITRI),
             NEIGH(1,ITRI),NSIDE(1,ITRI),IPROP(1,ITRI),
             NEIGH(2,ITRI),NSIDE(2,ITRI),IPROP(2,ITRI),
             NEIGH(3,ITRI),NSIDE(3,ITRI),IPROP(3,ITRI)
         ENDDO
       ENDIF
    ELSEIF (INDGRD(1).EQ.6) THEN
      IF (NLSLB.OR.NLCRC.OR.NLELL.OR.NLTRI) THEN
        RIA  RGA  RAA
      ELSEIF (NLPLG.OR.NLFEM.OR.NLTET) THEN
        XPCOR  YPCOR  ZPCOR
      ENDIF
    ENDIF
 ENDIF
 ** 2B. y- or poloidal grid surfaces
 NLPOL
 NLPLY  NLPLA  NLPLP
 NP2ND  NPSEP  NPPLA  NPPER
 IF (INDGRD(2).LE.5) THEN
    YIA  YGA  YAA  YYA
 ENDIF
 ** 2C. z- or toroidal grid surfaces
 NLTOR
 NLTRZ  NLTRA  NLTRT
 NT3RD  NTSEP  NTTRA  NTPER
 IF (INDGRD(3).LE.5) THEN
    ZIA  ZGA  ZAA  ZZA  ROA
 ENDIF
 ** 2D. mesh multiplication
 NLMLT
 NBMLT
 IF (NBMLT GT 1) (VOLCOR(NM), NM = 1,NBMLT)
 ** 2E. additional cells outside standard mesh
 NLADD
 NRADD
 IF (NRADD.GT.0) (VOLADD(NM), NM = 1,NRADD)
\end{verbatim}

\medskip

{\bf Meaning of the Input Variables}

\medskip

\begin{list}{ }{}
\item[{\bf INDGRD }]
This index controls the meaning of input variables of different
standard grid options.

INDGRD(1)  for the radial (or x-direction) grid

INDGRD(2)  for the poloidal (or y-direction) grid

INDGRD(3)  for the toroidal (or z-direction) grid
\end{list}

{\sl Data for first standard mesh: radial or x grid RSURF}

\begin{list}{ }{}
\item[{\bf NLRAD= .TRUE.}] ~

A radial or x grid is defined. Otherwise the complete
sub-block 2A may be omitted.
Depending upon the logical parameters in the next input card
the ``geometry - level" variable LEVGEO is set internally.
\begin{list}{ }{}
\item[{LEVGEO=1}]
cartesian coordinates x (and y)
\item[{LEVGEO=2}]
polar coordinates r (and $\theta$)
\item[{LEVGEO=3}]
general curvilinear coordinates: a full 2D mesh (polygonal
coordinate lines) is used in the x - y  plane. Grid cuts are
permitted in the y-direction.
\item[{LEVGEO=4}]
a 2D ``finite element" mesh of triangles is used in the x - y plane
\item[{LEVGEO=5}]
a 3D ``finite volume" mesh of tetrahedrons used
\item[{LEVGEO=10}]
a general, user defined geometry block is used.
All geometrical calculations are performed in problem specific routine
VOLUSR, TIMUSR, ...etc.
\end{list}
If NLRAD=.FALSE., then no spatial grid is defined
and the default geometry level LEVGEO = 1 is used.
Volume discretisation may still be achieved ``by hand" by defining ``additional
surfaces" (input block 3b) and appropriate cell number switching.

\item[{\bf INDGRD(1) }]
\begin{list}{ }{}
\item[{= 1 }]
standard grid option; the input parameters are used as described
below.
\item[{= 2 }]
As for INDGRD(1)=1, but
in case of LEVGEO = 2 by this option a radial grid is defined such
that the spacing
of radial surfaces is not equidistant at some poloidal position, but
instead such that the area enclosed between
two neighboring surfaces is kept constant. The input variable RGA
is irrelevant in this case.
\item[{= 3,4,5 }]
not in use
\item[{= 6 }]
Data for NR1ST radial surfaces
are read from the work array RWK. These
data must have been written onto RWK in the user supplied subroutine INFCOP.
By this option grid parameters may directly be transferred into
EIRENE from other files, e.g. from plasma transport codes
(see below: section \ref{sec2.14}. Data for interfacing routine INFCOP).
\end{list}
\end{list}

One and only one of the next 8 logical variables NLSLB $\dots$ NLTET, NLGEN must be .TRUE.
This variable defines the so called ``geometry level" of a run, i.e. the internal EIRENE variable LEVGEO

\begin{list}{ }{}
\item[{\bf NLSLB = .TRUE. }] ~

Geometry level: LEVGEO = 1

Cartesian geometry, the x co-ordinate is discretized by setting

$x_I$=RSURF(I) I=1,NR1ST .

Furthermore,
the flux-surface labeling grid RHOSRF(I) is identical with
the grid RSURF(I).
\item[{\bf NLCRC = .TRUE. }] ~

Geometry level: LEVGEO = 2

Cylindrical or toroidal geometry,
1D (``radial") mesh of concentric, circular surfaces. Polar coordinates are
used in the x-y plane. The third coordinate (either $z$ or toroidal angle $\phi$)
is either ignorable or discretized according to options described below (block 2c).

The radial surfaces are given by
$r^2 = x^2 + y^2 = const.$ and
radial coordinate $r$ is discretized by setting $r_I$=RSURF(I) I=1,NR1ST.
Furthermore RHOSRF(I) = $\sqrt{ \frac{AREA}{\pi} }$ where
AREA is the area inside surface
number $I$.
Thus for this option one has again: RHOSRF(I) = RSURF(I), I=1,NR1ST.
\item[{\bf NLELL = .TRUE. }] ~

Geometry level: LEVGEO = 2

Mesh of nested, but not necessarily concentric or confocal elliptical
flux surfaces.
The equation for the ``radial" surface
is $(x - EP)^2 + (  y/EL )^2 = r^2$.
The radial coordinate $r$ is discretized by setting
$r_I$=RSURF(I) I=1,NR1ST.

EP and EL may vary with coordinate $r$. These parameters are stored in the
arrays EP(I),EL(I), I=1,NR1ST which now are used in addition to
RSURF to define one coordinate surface in the first (radial or $x$- grid).

RHOSRF: as in NLCRC option.

Note: RHOSRF and RSURF may differ in this case.
\item[{\bf NLTRI = .TRUE. }] ~

Geometry level: LEVGEO = 2

to be written: triangularity in mesh of nested closed algebraic
surfaces
\item[{\bf NLPLG = .TRUE. }] ~

Geometry level: LEVGEO = 3

The mesh in the x-y plane
is described by NR1ST polygonal arcs of length NRPLG
each. A polygon
may consist of several ``valid" and ``invalid"
parts (to account for ``grid cuts" in CFD meshes).
The ``invalid" parts of a polygon are not seen by test particles and are
allowed for in EIRENE only in order to facilitate index mapping
in case of runs coupled to plasma transport models,
which resort to computer
generated meshes including grid cuts.

The polygons must not intersect each other.

In this case RHOSRF(1)=0., and RHOSRF(I) is the area enclosed by
polygon number 1 and polygon number I.
\item[{\bf NLFEM = .TRUE. }] ~

Geometry level: LEVGEO = 4

The mesh in the x-y plane
consists of NR1ST triangles, composed from NRKNOT knots.

In this case a flux surface labeling grid RHOSRF is not defined.
\item[{\bf NLTET = .TRUE. }] ~

Geometry level: LEVGEO = 5

3D discretisation of volume by tetrahedrons. For this grid option
please make contact to the authors.
\item[{\bf NLGEN = .TRUE. }] ~

Geometry level: LEVGEO = 10

Arbitrary geometrical configuration. Mesh consists of NR1ST arbitrarily
shapes cells (in any dimension).
Particle tracing routines
must be provided by user (VOLUSR, SAMUSR, TIMUSR, LEAUSR)
\medskip
\item[{\bf NR1ST }]
Number of grid-points in the radial (or x-direction) standard mesh

if NR1ST $\le$ 1, no radial (or x-direction)
standard mesh is defined.
\item[{\bf NRSEP }]
This flag is active for LEVGEO = 1 or LEVGEO = 2. Otherwise it is irrelevant.

The first (x- or radial) standard mesh is composed by two equidistant x- or radial grids
of co-ordinate surfaces with different grid density.
There are NR1ST-NRSEP+1 grid-points in the first, and NRSEP
grid-points in the second part. The grid-point RSURF(NR1ST-NRSEP+1)
belongs to both parts.
\item[{\bf RIA }]
left endpoint of standard grid (internally set $\ge 0$ if LEVGEO =2);
RSURF(1)=RIA
\item[{\bf RGA }]
boundary separating first and second part of standard grid with
different grid-point densities; RSURF(NR1ST-NRSEP+1)=RGA
\item[{\bf RAA }]
right endpoint of standard grid: RSURF(NR1ST)=RAA
\item[{\bf RRA }]
if RRA $>$ RAA, one additional, outer void zone is defined

RSURF(NR1ST)=RRA, and the parameter RAA now determines the surface no. NR1ST-1.

(irrelevant, if RRA $\le$ RAA)
\end{list}

if NLELL = .TRUE. :
\begin{list}{ }{}
\item[{\bf EPIN }]
Value of EP(r) for cylindrical co-ordinate surface number 1 with
$r_1$=RIA
(see: NLCRC = .TRUE. option above)
\item[{\bf EPOT }]
Value of EP(r) for cylindrical surface number NR1ST with

$r_{NR1ST}$=RAA
\item[{\bf EPCH }]
Value of EP(r) for cylindrical surface number NR1ST+1 with

$r_{NR1ST+1}$=RRA

(irrelevant if RRA $\le$ RAA)
\item[{\bf EXEP }]
The variation of the ``shift function" EP(r) with r is given by

$EP(r) = EPIN + { ( \frac{r - RIA}{RAA -
RIA} ) }^{EXEP}
\cdot (EPOT - EPIN)$
\item[{\bf ELIN }]
Value of EL(r) for cylindrical surface number 1 with \\
$r_1$=RIA (see: NLELL = .TRUE. option)
\item[{\bf ELOT }]
Value of EL(r) for cylindrical surface number NR1ST \\
with $r_{NR1ST}$=RAA
\item[{\bf ELCH }]
Value of EL(r) for cylindrical surface number NR1ST+1 with

$r_{NR1ST+1}$=RRA

(irrelevant for RRA $\le$ RAA)
\item[{\bf EXEL }]
The variation of the ``ellipticity function" EL(r) with
r is given by

$EL(r) = ELIN + { ( \frac{r - RIA}{RAA -
RIA} ) }^{EXEL}
\cdot (ELOT - ELIN)$
\end{list}

if NLTRI = .TRUE. :
\begin{list}{ }{}
\item[{\bf TRIN, TROT, TRCH, EXTR }]
to be written: the option for algbraically given triangular grids is currently
no available.

\end{list}
if NLPLG = .TRUE. :
\begin{list}{ }{}
\item[{\bf NR1ST }]
number of polygons for discretisation in ``radial" or x direction
\item[{\bf NRPLG }]
Number of points per polygon
\item[{\bf NPPLG }]
Number of valid parts on each polygon.
Each polygon is described by the x and y
co-ordinates of
NRPLG points. It is not necessary that all this
points are used for the
polygon. One can cut the polygon into several valid parts interrupted
by parts which are not seen by the test particles. ( Default :
NPPLG = 1 ). This option facilitates the use of 2-d computer generated
meshes which contain topological grid cuts.
\item[{\bf XPCOR,YPCOR }] ~

shift whole mesh by that vector in x,y-plane
\item[{\bf RFPOL }]
if RFPOL $>$ 0., one additional polygon zone is defined,
at a distance RFPOL outside the polygon NR1ST,
and then NR1ST is increased by one.

(irrelevant, if RFPOL $\le 0$.)
\item[{\bf NPOINT(1,J) }] ~

Index of the first point of the valid part number J
(same for each radial polygon)

( Default : NPOINT(1,1) = 1 )
\item[{\bf NPOINT(2,J) }] ~

Index of the last point of the valid part number J
(same for each radial polygon)

( Default : NPOINT(2,1) = NRPLG )
\item[{\bf XPOL(K,I) }]
x-co-ordinate of the polygon point number K on polygon number I
\item[{\bf YPOL(K,I) }]
y-co-ordinate of the polygon point number K on polygon number I
\end{list}
if NLFEM = .TRUE. :
\begin{list}{ }{}
\item[{\bf NR1ST }]
number of triangles for discretisation in x-y plane
\item[{\bf XPCOR,YPCOR }] ~

shift whole mesh by that vector in x,y-plane
\item[{\bf NRKNOT }]
There are NRKNOT knots, by which the triangles are defined
\item[{\bf XTRIAN,YTRIAN }] ~

x and y co-ordinates of the knots, respectively.
\item[{\bf NVERT(I,ITRI) }] ~

Each triangle ITRI is defined by 3 points $P_1, P_2$ and $P_3$ from
the set of NRKNOT knots.

NVERT(I,ITRI) is the number of
point $P_I$ (I=1,2,3) in the set of knots for triangle ITRI.
\item[{\bf NEIGH(I,ITRI) }] ~

The three sides of each triangle $S_1, S_2, S_3$ are defined such that
$S_1$ connects $P_1$ and $P_2$, $S_2$ connects $P_2$ and $P_3$ and $S_3$
connects $P_1$ and $P_3$.

NEIGH(I,ITRI) is the number of the neighboring triangle at side $S_I$
for I=1,2,3.
\item[{\bf NSIDE(I,ITRI) }] ~

NSIDE(I,ITRI) is the number (1,2 or 3) of the side of the neighboring
triangle, which corresponds to side $S_I$ of the triangle ITRI
to be written
\item[{\bf IPROP(I,ITRI) }] ~

ISTS = ABS(IPROP) is the integer, by which a
particular surface property is assigned to
side I of the triangle ITRI. ISTS=0 stands for default grid option (transparent
surface, cell indexing is done automatically). Otherwise ISTS is the number of
an additional (ISTS $<$ NLIMI) or non-default standard (NLIM $<$ ISTS $<$NLIM+NSTSI)
surface option as read in sub-blocks 3a or 3b, respectively.
By default the normal vector of each side of a triangle points out of the triangle.
In case IPROP $<$ 0, this vector points into the triangle.
This is relevant at surface options with ILSIDE $\not=$ 0.
\end{list}
if NLTET = .TRUE. :

Documentation not available here. Make contact with the authors.

if NLGEN = .TRUE. :
\begin{list}{ }{}
\item[{\bf NR1ST }]

number of cells in otherwise arbitrary mesh option NLGEN.
Use problem specific routines (see section \ref{sec3})
to set up mesh, cell volumes, and flight interesction times for test partilces.
\end{list}

{\sl Data for second standard mesh: poloidal or y grid PSURF}
\begin{list}{ }{}
\item[{\bf NLPOL }]
A poloidal or y grid is defined. Otherwise the complete block 2B
may be omitted and the volume averaged tallies are then automatically
integrated over this co-ordinate.
\item[{\bf INDGRD(2) = 1 }] ~

standard grid option
\item[{\bf INDGRD(2) = 2,3,4,5,6 }] ~

not in use  (default: INDGRD(2)=1)
\item[{\bf NP2ND }]
Number of grid-points in y- or poloidal direction
\item[{\bf YIA,YGA,YAA,YYA }] ~

\begin{list}{ }{}
\item[ for the LEVGEO=1 option: ] ~

the y grid PSURF is defined in the same way as the x grid, using the
parameters YIA,YGA,.... (cm) instead of RIA,RGA,....

\item[ for the LEVGEO=2 option: ] ~

the poloidal angle $\theta$ grid PSURF is defined in the same way as the radial $r$ grid was,
using
the parameters YIA,YGA,...(poloidal angle in degree) instead of RIA,RGA,...

for all options LEVGEO $>$ 2:
this input card is irrelevant.
\end{list}
Defaults: (needed for scaling, i.e., cell volumes):

YIA=0., YGA=0., YAA=1., YYA=1. in case of LEVGEO=1,

and YIA=0., YGA=0., YAA=360., YYA=360. in case of LEVGEO=2.

\end{list}
\medskip
{\sl Data for third standard mesh: toroidal or z grid TSURF}
\begin{list}{ }{}
\item[{\bf NLTOR }]
A toroidal or z grid is defined. Otherwise the complete block 2C
may be omitted, defaults described below are used and the volume averaged tallies are then automatically
integrated over this co-ordinate.

In case NLTOR = TRUE, sub-block 2C must be read.
\item[{\bf INDGRD(3) = 1 }] ~

standard grid option
\item[{\bf INDGRD(3) = 2,3,4,5,6 }] ~

not in use (default: INDGRD(3)=1)
\end{list}
One and only one of the next four boolean variables must be TRUE
\begin{list}{ }{}
\item[{\bf NLTRZ=TRUE }] ~

cylindrical approximation is used, i.e., TSURF is a grid in z
direction. The co-ordinate surfaces are given by z=TSURF(L)

(Default: NLTRZ = TRUE)

\item[{\bf NLTRA = TRUE }] ~

toroidal approximation is used. The coordinate line is
a polygonal approximation of a circle, or of an angular section thereof.

In case NLTOR = TRUE, the 3rd grid TSURF is a grid of
toroidal angles.
The toroidal segment (or the full torus) is approximated by NT3RD-1 straight cylindrical segments.

In case NLTOR = FALSE, there are NTTRA toroidal periodicity boundaries,
such that the toroidal segment is approximated, again, by NTTRAM=NTTRA-1 straight
cylinders, but without toroidal resolution in the results.
(Note: If NLTOR, then: NTTRA = NT3RD, internally.)

The torus axis of the entire mesh can be shifted in radial
direction by adding a radial offset ROA (see below) to the x
coordinates of the poloidal mesh (RSURF,PSURF).

The radial shift of the poloidal mesh
RMTOR of the approximated torus is computed from ROA
such that the volume
inside radial surface NR1ST is exactly equal to the volume
of an exact torus with poloidal cross section defined by the
shifted radial grid (first standard mesh: RSURF).

Due to the approximations made by
defining a torus by NTTRAM = NTTRA - 1 straight cylinders,
this condition is fulfilled only approximately
for the other radial surfaces. RMTOR converges to ROA with increasing
NTTRA. NTTRA $\simeq$ 30 is already a very good approximation.

(Default: NLTRA = FALSE).
\item[{\bf NLTRT = TRUE }] ~

torus co-ordinates R,PHI,THETA. Presently being developed for
NLSLB,NLPLG and NLTRI options.
Not ready for use.
\item[{\bf NT3RD }]
Number of grid-points in z- or toroidal direction

(default: NT3RD = 1, i.e. no grid is defined)
\item[{\bf NTTRA }]
only needed in case NLTRA and .NOT.NLTOR. See above.
\item[{\bf ZIA,ZGA,ZAA,ZZA,ROA }] ~

The 3rd  grid ZSURF is defined in the same way as the x grid, using
the parameters ZIA, ZGA,.... (cm) instead of RIA, RGA,....

In case of NLTRZ = TRUE , a z-grid is defined. ROA is irrelevant.

In case of NLTRA = TRUE ,
ZIA and ZAA are toroidal angles (in degrees). A grid of toroidal angles
is defined. For example use ZIA=0.0 and ZAA=360.0 for a full torus.
In case NLTOR this grid also defines the toroidal resolution.
If .NOT.NLTOR, then periodicity at the endpoints ZIA and ZAA is automatically
enforced.

ROA is the radial shift of the poloidal cross section defined by the x-y grids.
E.g.: if the x-y- grids are given with magnetic axis as origin, then ROA is the major radius.
If the x-y- grids are already given with respect to the torus axis at their origin, then ROA =0 (or better e.g.: $1.0E-4 = \mbox{ROA} \ll 1$)

Note: ROA affects evaluation of cell volumes.

Note also: in case of geometry and trajectory plots (input block 11),  this major radius offset ROA of poloidal cross sections has to be taken into
account when defining plot-frames.

(Defaults: NLTRA = FALSE, ZIA = 0 , ZAA = 1)

\end{list}
{\sl Data for mesh multiplication}
\begin{list}{ }{}
\item[{\bf NLMLT }]
the complete ``Standard Mesh" data are copied NBMLT times
\item[{\bf NBMLT }]
Number of identical copies of the standard mesh. Transition from one
such mesh (called ``block" in EIRENE) into another one has to be defined
by transparent additional surfaces (see block 3B)
\item[{\bf VOLCOR }]
The volumes of all cells of the standard mesh
as computed by EIRENE (in subroutine. VOLUME) are multiplied
by one common factor VOLCOR for each ``block".
\end{list}
{\sl Data for additional zones outside the Standard mesh}
\begin{list}{ }{}
\item[{\bf NLADD }]
There are cells in the computational volume, which are defined
through ``additional surfaces" as cell boundaries.
E.g. a standard mesh (if there is one) is augmented by ``additional
cells" in this case.
If NLADD = FALSE , the complete block 2D may be omitted and
the defaults are used.
\item[{\bf NRADD }]
Number of additional zones. The  cell indexing along test
flights in these zones has to be specified explicitly by making use of
the ILSWCH, ILCELL parameters (block 3B)

(Default: NRADD = 0 ).
\item[{\bf VOLADD }]
Volume ($cm^{-3}$) of each additional zone as seen by the
test-particles.
\end{list}
\subsection{Mesh Parameters} \label{sec2.2.1}
As stated above,
the computational volume in an EIRENE run is subdivided into NSBOX cells,

NSBOX = (NR1ST $\cdot$ NP2ND $\cdot$ NT3RD) $\cdot$ NBMLT + NRADD .

The can mesh consist two different parts.

The first is the so called ``standard mesh" part (``VOXEL-geometry"), defined
by the grids RSURF, PSURF and TSURF. There may then be NBMLT copies of that
mesh, separated by additional surfaces (see below, section \ref{sec2.3b}
(default: NBMLT=1).

The second part is the so called ``additional cell region". Up to
NRADD additional cells, of arbitrary shape, are defined by their
boundaries (``BREP"-geometry): the ``additional surfaces" defined in section
\ref{sec2.3b}.

\medskip

NR1ST is the number of surfaces in the 1st (radial -or x-) grid.
There are NR1STM = NR1ST - 1 cells.

The cell index of a particular cell is called NRCELL.

1$\le$ NRCELL $\le$ NR1STM

NP2ND is the number of surfaces in the 2ND (poloidal -or y-) grid.
There are NP2NDM = NP2ND - 1 cells.

The cell index of a particular cell is called NPCELL.

1$\le$ NPCELL $\le$ NP2NDM

NT3RD is the number of surfaces in the 2RD (toroidal -or z-) grid.
There are NT3RDM = NT3RD - 1 cells.

The cell index of a particular cell is called NTCELL.

1$\le$ NTCELL $\le$ NT3RDM

For the cells of the standard mesh, NACELL = 0

For the cells in the ``additional cell region" :

1$\le$ NACELL $\le$ NRADD

and

NRCELL = 0, NPCELL = 1, NTCELL = 1

\medskip

The last radial (or x-) cell (number NR1ST), for any NPCELL, NTCELL
inside the standard mesh,
is not a real geometrical cell. It is used to store
the radial (or x-) average of the volume averaged tallies.

Likewise, the last poloidal (or y-) cell (number NP2ND) is not a
real geometrical cell, but used to store
the poloidal (or y-) average of the volume averaged tallies (for any
NRCELL, NTCELL inside the standard mesh)

NTCELL = NT3RD: analogically.

\medskip

A particular cell can be specified in one of two possible ways:

Either by giving the 5 cell numbers:

NRCELL, NPCELL, NTCELL, NBLOCK, NACELL

or, alternatively, by giving the position in the 1-dimensional cell array

NCELL

\medskip
The relation between these two options is:

NCELL = NRCELL + ((NPECLL-1)*NR1P2) + (NTCELL-1))*NP2T3 + NBLCKA

with

NBLCKA = (NBLOCK -1) * NSTRDT +NRADD

\subsection{geometry level LEVGEO: symmetry and dimensionality of a run}
\label{sec2.2.2}

tbd: add a few figures here....



The discretisation of a domain is achieved by defining surfaces.
These can be sets of coordinate surfaces $u^i(x,y,z) = const_i$
(then producing a discretisation of the coordinate $\nabla u$).


{\bf LEVGEO = 1}
Simple cartesian grids (1D, 2D or 3D) are used for code development, comparison with simple (semi-analytical) transport models,
general 1D  ``one speed transport theory",
validation of diffusion approximations, etc......  The co-ordinates in the 3 directions are referred to as x,y and either z or phi, respectively,
the latter depending on whether the third co-ordinate is linear (cartesian) or circular.

{\bf LEVGEO = 2} Typically used in connection with 1D radial core
transport models (e.g., ``Duechs code", 1982 TRANSP 1987, RITM,
(1996),...), circular limiter tokamaks, or for simple parametric
studies in toroidal or cylindrical coordinates (this is the oldest
EIRENE geometry option, available since 1980). Shifted elliptical grid surfaces (1D),
but also in combination with a polar grid (then: 2D) and/or a cylindrical or toroidal
resolution in the third dimension.

{\bf LEVGEO = 3} This 2D discretisation of a domain, based upon polygonal lines,  in the x-y
plane was originally developed for the B2-EIRENE interfacing. More generally, this option allows to
directly use many 2D CFD meshes directly in EIRENE for kinetic transport.

{\bf LEVGEO = 4} This example shows the extension of the 2D
B2-grid, polynomial option LEVGEO = 3, into the outer
``vacuum"-region by utilizing the LEVGEO = 4 option (plenum, vacuum
chamber, grid refinement)

{\bf LEVGEO = 5} general 3D grid based upon tetrahedrons

{\bf LEVGEO = 10} user supplied geometry option. All grid relevant computations are
done in user specified (problem specific) routines: EIRENE module:  USER.f
